% 通信仿真链路（由 body/chap03.tex 的 tikzpicture 以 \input 引入，单位 cm）
% 走线原则：发送链在上排自左向右，经信道模型（虚线框）落到下排，接收链在下排自右向左回折，
% 使发送与接收的箭头方向相反，避免两排同向造成误读；指标统计置于底部，
% 由接收判决与发送比特真值（点划线）两路输入；蒙特卡洛控制沿最左侧竖直虚线回到帧生成。
% 各节点用 text width 固定文字宽度，防止中文撑宽后与相邻节点边框贴合。
\tikzset{
  simblk/.style={draw, rectangle, rounded corners=2pt, minimum height=0.64cm,
    inner sep=2pt, align=center, font=\footnotesize, line width=0.5pt},
  simch/.style={draw, rectangle, minimum height=0.58cm,
    inner sep=2pt, align=center, font=\footnotesize, line width=0.5pt, fill=black!4},
  simsum/.style={draw, circle, inner sep=0pt, minimum size=0.38cm,
    font=\footnotesize, line width=0.5pt},
  simarr/.style={-latex, line width=0.5pt},
  simtruth/.style={-latex, line width=0.45pt,
    dash pattern=on 4pt off 2pt on 1pt off 2pt},
  simfb/.style={-latex, line width=0.45pt, dashed},
  simtag/.style={font=\footnotesize}
}

% ---------------- 发送链（上排，自左向右） ----------------
\node[simblk, text width=1.60cm] (t-gen) at (0.95,3.50) {帧生成};
\node[simblk, text width=2.60cm] (t-mod) at (3.69,3.50) {SFH-MFSK 调制};
\node[simblk, text width=1.60cm] (t-hop) at (6.43,3.50) {跳频混频};
\draw[simarr] (t-gen.east) -- (t-mod.west);
\draw[simarr] (t-mod.east) -- (t-hop.west);

% ---------------- 信道模型（虚线框内：多径、时缩、噪声与干扰） ----------------
\draw[dashed, line width=0.45pt] (7.90,2.90) rectangle (14.60,0.05);
\node[simtag, anchor=east] at (14.48,2.66) {信道模型};
\node[simch, text width=2.60cm] (c-mp) at (9.40,2.20) {多径 $h(\tau,t)$};
\node[simch, text width=2.60cm] (c-dop) at (9.40,1.30) {多普勒时缩};
\node[simsum] (c-sum) at (9.40,0.50) {$+$};
\node[simch, text width=3.40cm] (c-noise) at (12.75,0.50) {噪声 $n(t)$ / 干扰 $J(t)$};
\draw[simarr] (t-hop.east) -- (9.40,3.50) -- (c-mp.north);
\node[simtag, anchor=south] at (7.80,3.56) {$s(t)$};
\draw[simarr] (c-mp.south) -- (c-dop.north);
\draw[simarr] (c-dop.south) -- (c-sum.north);
\draw[simarr] (c-noise.west) -- (c-sum.east);

% ---------------- 接收链（下排，自右向左回折） ----------------
\node[simblk, text width=1.60cm] (r-fe) at (13.55,-1.30) {接收前端};
\node[simblk, text width=3.60cm] (r-syn) at (10.06,-1.30) {同步（重复 LFM + 图案）};
\node[simblk, text width=2.10cm] (r-dem) at (6.62,-1.30) {Notch 解调};
\node[simblk, text width=2.40cm] (r-dec) at (3.63,-1.30) {判决与帧校验};
\draw[simarr] (c-sum.south) -- (9.40,-0.62) -- (13.55,-0.62) -- (r-fe.north);
\node[simtag, anchor=south] at (11.60,-0.56) {$r(t)$};
\draw[simarr] (r-fe.west) -- (r-syn.east);
\draw[simarr] (r-syn.west) -- (r-dem.east);
\draw[simarr] (r-dem.west) -- (r-dec.east);

% 接口量标注：置于相邻两级间隙下方，避免压住节点边框
\node[simtag, anchor=north] at (12.30,-1.70) {$\hat\tau,\hat k_0$};
\node[simtag, anchor=north] at (7.97,-1.70) {$\hat m_k$};
\node[simtag, anchor=north] at (5.20,-1.70) {判决比特};

% ---------------- 指标统计 ----------------
\node[simblk, text width=8.10cm, minimum height=1.02cm] (s-metric) at (6.60,-2.85)
  {指标统计\\[1pt]$P_b$、$P_{\mathrm{syn}}$、$P_{\mathrm{frame}}$、运算量};
\draw[simarr] (r-dec.south) -- (3.63,-2.33);

% 发送比特真值：由帧生成直下，与判决比特在指标块内比对
\draw[simtruth] (t-gen.south) -- (0.95,-2.85) -- (s-metric.west);
\node[simtag, rotate=90, anchor=south] at (0.58,-0.95) {发送比特（真值）};

% ---------------- 蒙特卡洛控制 ----------------
\draw[simfb] (s-metric.south) -- (6.60,-3.90) -- (-0.55,-3.90) -- (-0.55,3.50) -- (t-gen.west);
\node[simtag, anchor=north] at (3.20,-3.96) {蒙特卡洛重复 $N_{\mathrm{fr}}$ 帧};




